La modernisation d’une électronique vieillissante et critique pour la sécurité est souvent plus rentable et plus durable sur le plan environnemental que son remplacement complet, à condition que l’équivalence fonctionnelle avec le système d’origine puisse être démontrée et le système modernisé certifié à nouveau selon des normes actuelles. Cette étude de cas, qui s’inscrit dans la continuité de travaux antérieurs, aborde ce défi à travers l’exemple d’un module existant de test et de diagnostic d’un système anti-enrayeur, utilisé dans le transport ferroviaire urbain européen.

Une carte contrôleur 8 bits existante, basée sur l’Intel~MCS-48, a été mise à niveau au moyen d’un circuit intégré prédiffusé programmable (FPGA) \texttt{Actel~A3P1000} à mémoire flash, qui émule l’unité centrale de traitement (CPU) et la logique d’entrées-sorties (E/S) d’origine, complété par un microcontrôleur (µC) Arm\textsuperscript{\textregistered} 32 bits isolé galvaniquement, un \texttt{STM32F401RET6}, pour une journalisation étendue des données. Un objectif central était l’élaboration d’un cadre méthodologique de test se prêtant à une application plus large dans la mise à niveau de systèmes critiques pour la sécurité. Au cœur de ce cadre se trouve une stratégie de vérification en deux couches, dérivée du modèle en~V : (1) Essais des composants et des interfaces : la simulation déterministe en langage de description de matériel (HDL), les contrôles JTAG et les mesures temporelles sur le matériel examinent l’exactitude logique, l’intégrité de base du matériel et le comportement temporel des interfaces (par exemple chien de garde, balayage en fréquence). (2) Essais au niveau système : des essais fondés sur des scénarios, menés sur un banc d’essai dédié~\cite{dolata2023} avec des cas de test définis en YAML, examinent le comportement intégré (mémoire des défauts, autotests, diagnostic). La simulation et les essais intégrés fondés sur des scénarios indiquent l’équivalence fonctionnelle du traitement des défauts par le cœur critique pour la sécurité (détection, mémorisation, affichage et lecture des défauts). Le balayage en fréquence s’écarte toutefois de l’original, et cet écart reste ouvert : dans la seule mesure comparative avec la carte d’origine, le balayage de l’émulation se termine environ 70~Hz plus bas (environ 850~Hz contre environ 920~Hz) et, d’après ses valeurs de conception, se déroule plus vite que la variation de fréquence d’environ 180~ms que~\cite{dolata2023} indique pour la carte de test d’origine ; pour trancher, il faut une nouvelle mesure de la carte d’origine ou une vérification de son circuit de balayage. Deux modifications du matériel changent par conception le comportement du module : les entrées STORE et RECALL du connecteur système n’agissent plus sur la mémoire des défauts, et une détection de sous-tension ajoutée peut réinitialiser le contrôleur émulé ; leur incidence dans le véhicule reste ouverte. Les éléments de preuve ne sont pas encore complets : les simulations et mesures archivées datent de 2023 et sont donc antérieures à la reconstruction ; sur l’appareil reconstruit, une équipe ayant poursuivi les travaux a exécuté la suite de tests du banc d’essai sans qu’un journal ait été archivé, et plusieurs scénarios, dont l’autotest, ne sont pas documentés pour cet appareil. Six défauts d’intégration documentés au cours de la reconstruction sont classés selon la manière dont ils ont été découverts ; ils n’ont pas été mis en évidence lors d’une application planifiée de la méthode en deux couches, et aucune méthode ne les a tous trouvés à elle seule. Les couches ont été mises en correspondance avec les normes IEC~61508~Ed.~2, EN~50129:2018, EN~50716:2023 et IEEE~1012:2024 dans leur ensemble, sans évaluation de la conformité ; EN~50155:2021 relève de travaux futurs. Bien que des installations d’essai complètes de type matériel dans la boucle (HIL) aient dépassé le cadre de ce projet, la démarche s’inspire de la pratique HIL afin de faciliter une intégration future.

La remise à neuf de l’électronique existante du matériel roulant peut réduire les coûts et éviter les émissions grises liées à la fabrication de nouveaux composants, et les mises à niveau à l’échelle de la carte électronique relèvent de la même logique. Une fonction anti-enrayeur (WSP) adaptative, destinée également aux flottes existantes et certifiée selon la norme EN~15595 en vigueur, constitue un précédent ; la validation de la méthodologie sur d’autres cas de mise à niveau reste à mener dans des travaux futurs.
